\documentclass[aspectratio=169,11pt]{beamer}

\usetheme{Madrid}
\usecolortheme{default}
\usefonttheme{professionalfonts}
\usepackage[T1]{fontenc}
\usepackage[utf8]{inputenc}
\usepackage{lmodern}
\usepackage{amsmath,amssymb}
\usepackage{booktabs}
\usepackage{ragged2e}
\usepackage{xcolor}
\usepackage{hyperref}

\definecolor{uiblue}{RGB}{38,63,106}
\definecolor{softblue}{RGB}{235,241,249}
\definecolor{warnred}{RGB}{155,38,38}
\setbeamercolor{structure}{fg=uiblue}
\setbeamercolor{frametitle}{fg=white,bg=uiblue}
\setbeamercolor{title}{fg=white,bg=uiblue}
\setbeamercolor{block title}{fg=white,bg=uiblue}
\setbeamercolor{block body}{bg=softblue}
\setbeamercolor{alerted text}{fg=warnred}
\setbeamertemplate{navigation symbols}{}
\setbeamertemplate{footline}[frame number]
\setbeamersize{text margin left=7mm,text margin right=7mm}

\title{Noen tanker}
\author{MHJ inputs}
\institute{UiO}
\date{September 13, 2026}

\newcommand{\good}{\textcolor{green!45!black}{\textbf{Strength}}}
\newcommand{\revise}{\textcolor{warnred}{\textbf{Revise}}}

\begin{document}

\begin{frame}[plain]
  \titlepage
\end{frame}

\begin{frame}{Great application}
\begin{block}{Overall impression}
A strong application with a distinctive central result and a potentially compelling independent program. 
\end{block}
\begin{itemize}
  \item Strong alignment with theoretical and computational quantum science at UiO.
  \item Strong and very credible  combination of condensed-matter theory, tensor networks and quantum technology.
  \item May be  an \alert{overextension}: several major methodological programmes are promised simultaneously?
  \item Some claims in Sections 2.2 and 2.5 may require require some clarifications/corrections
\end{itemize}
\end{frame}

\begin{frame}{Strong sides}
\begin{itemize}
  \item A memorable scientific starting point: kinetic frustration as a route to RVB physics.
  \item Clear continuity from the an excellent  publication record to an independent program.
  \item Exact few-hole states supply unusual benchmarks for both classical and quantum methods.
  \item Strong potential links across theoretical physics, condensed matter, materials research and UiO's quantum initiatives.
  \item Convincing teaching profile, program-level engagement and attention to student belonging.
  \item Explicit discussion of research-group development, external funding and computational resources.
\end{itemize}
\end{frame}

\begin{frame}{Petit corrections before submission}
\begin{enumerate}
  \item Remove the visible draft markers:
    \begin{itemize}
      \item ``ADD?'' in Section 2.4;
      \item ``Bond correlators in quantum gas microscopy?'' in Section 4.1;
      \item incomplete References 22 and 33.
    \end{itemize}
  \item Correct ``several orders of magnitudes'' to ``several orders of magnitude.''
  \item Correct ``frustrated systems finite doping'' to ``frustrated systems at finite doping.''
  \item Correct ``two- and three-dimension'' to ``two and three dimensions.''
  \item Standardise citations as, for example, \texttt{[8,9]} rather than \texttt{[8][9]}.
\end{enumerate}
\end{frame}

\begin{frame}{Can you check this? Finite $U$ stuff}
You write  that exchange frustration becomes more important as $U$ increases. Can you check this?  I may mess up things.

\begin{block}{Correct large-$U$ relation}
For the Hubbard model, the leading antiferromagnetic superexchange is
\[
J \simeq \frac{4t^2}{U},
\qquad
\frac{J}{t}\simeq\frac{4t}{U}.
\]
Thus, $J/t$ \emph{decreases} as $U/t$ increases. Have I misunderstood something here?
\end{block}

\textbf{Suggested replacement?:}

``Relaxing the infinite-$U$ constraint by reducing $U/t$ increases $J/t$ and therefore strengthens the role of exchange relative to kinetic processes. This connects the kinetic-frustration regime continuously to exchange-dominated frustrated magnetism.''
\end{frame}

\begin{frame}{Perhaps make a clearer hierarchy?}
The current plan contains at least six demanding directions:
\begin{columns}[T]
\column{0.48\textwidth}
\begin{itemize}
  \item finite-density iDMRG;
  \item fermionic neural quantum states;
  \item topological diagnostics;
\end{itemize}
\column{0.48\textwidth}
\begin{itemize}
  \item unconventional superconductivity;
  \item quantum algorithms and hardware;
  \item quantum Nematic Bond Theory.
\end{itemize}
\end{columns}

\begin{block}{Perhaps do the following?}
\begin{enumerate}
  \item \textbf{Scientific core:} survival and evolution of kinetic-frustration-induced RVB physics at finite density.
  \item \textbf{Principal methods:} tensor networks and validated NQS.
  \item \textbf{Exploratory branches:} quantum algorithms and qNBT, conditional on personnel, funding and preliminary results.
\end{enumerate}
\end{block}
\end{frame}

\begin{frame}{Scientific basis of the NQS programme}
\good
\begin{itemize}
  \item Variational Monte Carlo with NQS is a plausible route to frustrated fermionic systems in two and three dimensions.
  \item Determinant-based states, Jastrow correlators and neural backflow are relevant architectural ingredients.
  \item Exact few-hole RVB states provide stronger validation targets than energy comparisons alone.
  \item The square-lattice $t$--$J$ study supports the feasibility of hidden-fermion determinant states across several dopings.
\end{itemize}

\revise
\begin{itemize}
  \item Perhaps not  equate parameter efficiency with computational efficiency.
  \item Avoid saying that genuine 2D and 3D systems are simply ``beyond the reach'' of conventional methods.
  \item Specify system sizes and geometries when saying that the ``entire doping regime'' was accessed.
\end{itemize}
\end{frame}

\begin{frame}{NQS parameter counts: a more defensible claim}
\textbf{Current implication:} fewer parameters demonstrate that NQS outperform tensor networks.

\medskip
This comparison does not determine:
\begin{itemize}
  \item wall-clock cost or memory use;
  \item sampling and optimisation difficulty;
  \item reliability across random initialisations;
  \item accuracy of long-range observables;
  \item finite-size or thermodynamic-limit accuracy.
\end{itemize}

\begin{block}{Suggested wording}
``For several frustrated spin models, NQS have obtained energies competitive with leading tensor-network calculations while using substantially fewer explicit variational parameters. This parameter efficiency does not by itself imply lower total computational cost, and performance remains dependent on architecture, optimisation and sampling.''
\end{block}
\end{frame}

\begin{frame}{NQS: fermionic structure must be described precisely}
The roles of determinants, Jastrow factors and backflow may not be conflated

\begin{itemize}
  \item A Slater or hidden-fermion determinant can enforce antisymmetry.
  \item A symmetric Jastrow factor captures correlations but does not impose antisymmetry by itself.
  \item Neural backflow modifies orbitals or coordinates while retaining determinant antisymmetry.
\end{itemize}

\begin{block}{Suggested wording}
``Fermionic antisymmetry will be imposed through determinant-based reference structures, including hidden-fermion determinants. Symmetric Jastrow correlators and neural backflow transformations will then increase expressivity while preserving the required exchange symmetry.''
\end{block}

The plan could also specify how it imposes no double occupancy, fixed particle number, $S_z$ and lattice symmetries.
\end{frame}

\begin{frame}{NQS: proposed validation ladder}
\begin{enumerate}
  \item Exact diagonalisation on small clusters.
  \item Exact few-hole RVB states at increasing system size.
  \item Comparison with iDMRG on finite-width cylinders.
  \item Finite-density two-dimensional calculations.
  \item Pyrochlore applications only after the preceding stages succeed.
\end{enumerate}

\begin{block}{Quantitative validation criteria}
Relative energy error; energy variance; fidelity where available; spin, charge, dimer and hole correlations; structure factors; symmetry quantum numbers; multiple optimisation seeds; architecture ablations; and finite-size scaling.
\end{block}
Most likely, my bias here, 
accurate variational energy alone does not establish accurate long-range order, excitation physics or topology.
\end{frame}

\begin{frame}{NQS: entanglement and topology}
Rényi entropies have already been estimated for several NQS architectures, including autoregressive states. The open problem could therefore be stated more narrowly.

\begin{block}{Suggested wording}
``Rényi entropies can be estimated for several classes of NQS, but extending these estimators reliably to doped fermionic states and extracting subleading topological contributions in two dimensions remain challenging. I will test swap-operator and ratio-based estimators against exact and iDMRG benchmarks before applying them to larger systems.''
\end{block}

\begin{itemize}
  \item Flux insertion can reveal spectral flow, pumping and transitions between sectors.
  \item It does not, by itself, prove topological ground-state degeneracy.
  \item Finite-cylinder extraction of topological entanglement entropy requires careful width and boundary scaling.
\end{itemize}
\end{frame}

\begin{frame}{Promising future directions for NQS, here I asked ChatGPT/Claude}
\begin{itemize}
  \item Symmetry-equivariant and lattice-aware architectures.
  \item Autoregressive sampling to reduce Markov-chain limitations.
  \item Determinant or Pfaffian references with neural correlators.
  \item Neural backflow and hidden-fermion constructions.
  \item Transfer learning across doping, $J/t$ and system size.
  \item Hybrid NQS--tensor-network benchmarking.
  \item Excited-state and dynamical extensions.
  \item Ensembles or repeated seeds for optimisation uncertainty.
  \item Physics-informed networks encoding local dimer constraints.
  \item Supervised pretraining on exact RVB amplitudes, followed by continuation in doping or $J/t$.
\end{itemize}
\end{frame}

\begin{frame}{Section 2.5: separate three quantum tasks}
The section could  distinguish three levels of difficulty:

\begin{enumerate}
  \item \textbf{Prepare a known few-hole state.}\\
  Fidelity and observables can be certified classically.
  \item \textbf{Prepare an unknown finite-density ground state.}\\
  Success cannot be established from energy alone and may require costly optimisation or filtering.
  \item \textbf{Measure dynamical observables.}\\
  Time evolution, ancillas, sampling and noise introduce additional resource demands.
\end{enumerate}

\begin{block}{Possible framing}
Present this as a staged, benchmark-driven programme rather than implying near-term classical advantage.
\end{block}
\end{frame}

\begin{frame}{MPS preparation: state the necessary conditions}
I think that an  exact finite-state MPS representation may  not automatically imply efficient preparation.

\begin{itemize}
  \item Cost depends on system size, bond dimension, target fidelity and hardware connectivity.
  \item Mapping a 2D or 3D state to a 1D ordering can produce rapidly growing bond dimension.
  \item Standard sequential MPS circuits may be too deep for present hardware.
  \item Shallow constructions apply to special classes or rely on approximation, measurement and feedback.
\end{itemize}

\begin{block}{Suggested wording}
``For small clusters, the exact RVB states can be represented as MPS and compiled into quantum circuits. I will determine how bond dimension, two-qubit gate count, circuit depth and fidelity scale with system size, lattice ordering and hole number.''
\end{block}
\end{frame}

\begin{frame}{Quantum algorithms: missing resource specification}
The proposal could possibly specify the $t$--$J$ encoding and the resources to be measured.

\begin{columns}[T]
\column{0.48\textwidth}
\textbf{Encoding choices}
\begin{itemize}
  \item three-state local encoding;
  \item constrained qubit encoding;
  \item fermion-to-qubit mapping;
  \item fixed particle number and $S_z$;
  \item lattice connectivity and swaps.
\end{itemize}
\column{0.48\textwidth}
\textbf{Resource measures}
\begin{itemize}
  \item logical and physical qubits;
  \item two-qubit gate count and depth;
  \item ancillas and postselection rate;
  \item circuit repetitions;
  \item controlled-evolution cost;
  \item noise sensitivity.
\end{itemize}
\end{columns}
\end{frame}

\begin{frame}{Variational, adiabatic and Rodeo preparation}
\textbf{Variational methods:} address barren plateaus, noisy optimisation, symmetry leakage and limited shallow-circuit expressivity.

\medskip
\textbf{Adiabatic methods:} specify the interpolation Hamiltonian and test whether a small gap or phase transition obstructs preparation.

\medskip
\textbf{Resolution refinement plus Rodeo:}
\begin{itemize}
  \item define ``resolution refinement'' explicitly;
  \item report target-state overlap and filtering success probability;
  \item include energy resolution, evolution time, ancilla measurements and postselection cost;
  \item replace ``will make it perform better'' with a quantitative hypothesis.
\end{itemize}

\begin{block}{Testable formulation}
``I will test whether resolution-refined initial states increase the target-state overlap and success probability of Rodeo filtering, while measuring the resulting controlled-evolution depth and sampling cost.''
\end{block}
\end{frame}

\begin{frame}{Dynamical observables and spectral functions}
The dynamical spin structure factor requires more than measuring and Fourier transforming:
\begin{itemize}
  \item reliable real-time evolution;
  \item an unequal-time correlation protocol;
  \item sufficient evolution time for frequency resolution;
  \item management of decoherence and sampling uncertainty.
\end{itemize}

The cited ARPES-inspired algorithm estimates $A(k,\omega)$. It does not directly output a holon dispersion. In a fractionalised state, electron spectral weight may form a spin--charge continuum.

\begin{block}{Suggested wording}
``I will investigate whether momentum- and frequency-resolved features of the single-particle spectral function can provide signatures of holon dynamics and spin--charge separation, beginning with systems that can be verified by exact diagonalisation and classical time evolution.''
\end{block}
\end{frame}

\begin{frame}{Superconductivity and cross-validation}
\textbf{Holon condensation}
\begin{itemize}
  \item Here I am lacking knowledge....
\end{itemize}

\textbf{Cross-validation}
\begin{itemize}
  \item Agreement between qNBT and NQS increases confidence but may not prove correctness.
  \item Use ED, DMRG and QMC benchmarks wherever their respective regimes permit.
\end{itemize}
\end{frame}

\begin{frame}{Parhaps make an implementation matrix?}
\scriptsize
\setlength{\tabcolsep}{3pt}
\begin{tabular}{p{2.4cm}p{2.8cm}p{2.5cm}p{2.7cm}p{2.8cm}}
\toprule
\textbf{Work package} & \textbf{Main question} & \textbf{Method} & \textbf{First benchmark} & \textbf{Fallback outcome} \\
\midrule
Finite-density RVB & Does RVB physics persist? & iDMRG and NQS & Exact states and ED & Cylinder phase diagram \\
\addlinespace
NQS development & Can 2D/3D states be represented reliably? & Determinants, backflow, symmetry-aware NQS & Exact states and iDMRG & Validated architecture limits \\
\addlinespace
Quantum algorithms & Can exact RVB states be prepared economically? & MPS compilation, filtering, continuation & Exact fidelity & Resource scaling and small demonstration \\
\addlinespace
qNBT & Can NBT be extended to quantum spins? & Schwinger bosons and fluctuating fields & Square-lattice antiferromagnet & Controlled high-risk assessment \\
\bottomrule
\end{tabular}
\end{frame}

\begin{frame}{Think of stating explicitely}

\begin{itemize}
  \item which scientific ideas originated from you;
  \item which components extend existing collaborations;
  \item which methods you can deploy immediately;
  \item which expertise needs to be recruited or developed;
  \item what can you begin without external funding;
  \item what depends on PhD students or PD researchers.
\end{itemize}

\end{frame}


\begin{frame}{Final wordings?}
\begin{block}{Perhaps add something like this?}
``I possess a rare exact result that opens a well-defined new research direction, and I have a disciplined sequence of complementary methods with which to determine how far its physics survives.''
\end{block}
All in all
\begin{itemize}
  \item Strong application
  \item Scientific identity is distinctive and well connected to UiO.
  \item The key is to distinguish established results, testable hypotheses and longer-term possibilities.
\end{itemize}
\end{frame}


\end{document}
